\chapter{Gestión de Riesgos}
\label{sec:riesgos}

\section{Identificación y Categorización de Riesgos}
\label{sec:identificacion-riesgos}

Los riesgos asociados a la calidad de SpeechNotes se clasifican en dos categorías principales: \textbf{Riesgos del Producto} (asociados a la arquitectura, dependencias y comportamiento del software en producción) y \textbf{Riesgos del Proyecto de Pruebas} (asociados a la planificación, recursos y ejecución del sprint SQA).

\subsection{Riesgos del Producto (Software & Arquitectura)}
\label{subsec:riesgos-producto}

\begin{itemize}
    \item \textbf{RP-01: Caída o indisponibilidad de la API externa de NVIDIA NIM Cloud.}
    \begin{itemize}
        \item \textit{Probabilidad:} Media \quad \textit{Impacto:} Alto \quad \textit{Severidad:} Alta
        \item \textit{Plan de Mitigación:} Disponer de claves de API de respaldo y construir stubs/mocks desconectados en \texttt{pytest} para pruebas del backend sin dependencia externa.
    \end{itemize}
    \item \textbf{RP-02: Ausencia o falla en el dispositivo físico de entrada de audio (Micrófono).}
    \begin{itemize}
        \item \textit{Probabilidad:} Alta \quad \textit{Impacto:} Bajo \quad \textit{Severidad:} Media
        \item \textit{Plan de Mitigación:} Diseñar casos de prueba diferenciados (\texttt{CP-03}) que validen el manejo de errores por falta de permisos o hardware, complementando con simuladores de stream de audio en backend.
    \end{itemize}
    \item \textbf{RP-03: Inestabilidad o crasheos no controlados en la UI del frontend (React/Next.js).}
    \begin{itemize}
        \item \textit{Probabilidad:} Media \quad \textit{Impacto:} Medio \quad \textit{Severidad:} Media
        \item \textit{Plan de Mitigación:} Captura de logs de consola en Cypress y ejecución de pruebas exploratorias especificas para identificar pérdidas de estado en la UI.
    \end{itemize}
\end{itemize}

\subsection{Riesgos del Proyecto de Pruebas (Gestión & SQA)}
\label{subsec:riesgos-proyecto}

\begin{itemize}
    \item \textbf{RPJ-01: Tiempo insuficiente para completar automatización E2E al 100\%.}
    \begin{itemize}
        \item \textit{Probabilidad:} Alta \quad \textit{Impacto:} Medio \quad \textit{Severidad:} Alta
        \item \textit{Plan de Mitigación:} Adoptar una estrategia híbrida: priorizar la ejecución manual estructurada de los 14 casos de prueba críticos y automatizar únicamente los endpoints backend estables.
    \end{itemize}
    \item \textbf{RPJ-02: Discrepancias o problemas de compatibilidad en entornos locales de desarrollo.}
    \begin{itemize}
        \item \textit{Probabilidad:} Media \quad \textit{Impacto:} Alto \quad \textit{Severidad:} Alta
        \item \textit{Plan de Mitigación:} Utilizar la especificación de contenedores Docker (\texttt{docker-compose.yml}) y entornos virtuales unificados (\texttt{.venv}) para garantizar entornos replicables.
    \end{itemize}
    \item \textbf{RPJ-03: Documentación insuficiente de la API REST y eventos de WebSockets.}
    \begin{itemize}
        \item \textit{Probabilidad:} Baja \quad \textit{Impacto:} Alto \quad \textit{Severidad:} Media
        \item \textit{Plan de Mitigación:} Realizar sesiones de prueba exploratoria guiada y consultar directamente las especificaciones OpenAPI de FastAPI en \url{/docs}.
    \end{itemize}
\end{itemize}

\begin{longtable}{lllcp{4cm}}
\caption{Matriz consolidada de riesgos de Producto y Proyecto}
\label{tab:riesgos}\\
\toprule
\textbf{Código} & \textbf{Categoría} & \textbf{Prob / Imp} & \textbf{Severidad} & \textbf{Plan de Mitigación}\\
\midrule
\endfirsthead
\toprule
\textbf{Código} & \textbf{Categoría} & \textbf{Prob / Imp} & \textbf{Severidad} & \textbf{Plan de Mitigación}\\
\midrule
\endhead
\bottomrule
\endfoot

RP-01 & Producto & Media / Alto & Alta & Claves API respaldo y mocks en pytest\\

RP-02 & Producto & Alta / Bajo & Media & Casos de prueba sin micrófono (CP-03)\\

RP-03 & Producto & Media / Medio & Media & Captura de logs en Cypress e inspección\\

RPJ-01 & Proyecto & Alta / Medio & Alta & Estrategia híbrida (manual + pytest)\\

RPJ-02 & Proyecto & Media / Alto & Alta & Despliegue estandarizado con Docker\\

RPJ-03 & Proyecto & Baja / Alto & Media & Pruebas exploratorias sobre OpenAPI docs\\

\end{longtable}

\section{Resultados de Mitigación}
\label{sec:resultados-mitigacion}

\begin{itemize}
    \item \textbf{NVIDIA NIM (RP-01):} Se utilizó la clave de API activa de respaldo. Se completaron exitosamente las pruebas de formateo (\texttt{CP-10}) y chat contextual (\texttt{CP-11}).
    \item \textbf{Inestabilidad frontend (RP-03):} Se identificó y aisló el defecto crítico \texttt{BUG-003} (crasheo del chat sin transcripción seleccionada), logrando su corrección antes del cierre del sprint.
    \item \textbf{Tiempo de automatización (RPJ-01):} Se garantizó el 100\% de ejecución manual de los 14 casos de prueba con soporte de scripts en \texttt{pytest}.
    \item \textbf{Micrófono (RP-02):} Se ejecutó \texttt{CP-03} (sin dispositivo) registrando el manejo correcto de alertas en la interfaz.
    \item \textbf{Manejo Global de Defectos:} Se registraron 8 defectos en Jira Software, logrando resolver el 100\% de los bloqueantes.
\end{itemize}
